En la figura es mostren tres fils conductors rectilinis i infinitament llargs, perpendiculars al pla del paper, per cadascun dels quals circula una mateixa intensitat de corrent de 0,30~A en el sentit que va cap a dins del paper. Aquests tres conductors estan situats en tres vèrtexs d'un quadrat de 0,20~m de costat.

\figura[0.25]{fig1.png}

\begin{apartat}{1}
Representeu en un esquema els camps magnètics, en el vèrtex $C$, generats pels conductors $A$ i $B$, i també el camp total. Calculeu el mòdul del camp magnètic total en aquest punt.
\end{apartat}

\begin{apartat}{1}
Representeu la força total sobre el conductor $C$ i calculeu el mòdul de la força que suporten 2,00~m del conductor que passa per $C$.
\end{apartat}

\nota{El mòdul del camp magnètic a una distància $r$ d'un fil infinit pel qual circula una intensitat $I$ és: $B = \dfrac{\mu_0 I}{2\pi r}$, on $\mu_0 = 4\pi\times 10^{-7}\ \mathrm{T\,m\,A^{-1}}$.}
